\scanpage{56}
% archive sheet number: 55
$\Sigma^*_{\underline J}$ défini \deleted{par} à partir de \deleted{la fibré} la cônique \deleted{\unsure{déf}} d'équation homogène
\[
 \sum_{i\in J}\unsure{y|x}_i^2=0 \tag*{(12)}
\]
% slanted
\margin{description valable loc\textsuperscript{t}. (ét)~; sinon, il faut prendre, dans $\check{\mathbb P}(p_*(\mathcal O_{\underline J}))$, la conique d'équation \deleted{\unsure{de}} $Q=0$, où $Q(u\deleted{\unsure{x|\otimes}})=\mathrm{Tr}\,u^2$ (NB $p:\underline J\longrightarrow S$ le morphisme structural)}
sur laquelle le groupe $V_{\underline J}$ \inserted{$\simeq\mathbb F_2^{\underline J}/\mathbb F_2$} opère \deleted{par un} \inserted{par} changement de signes sur les $x_i$, avec comme quotient $\Sigma_{\underline J}$, défini par l'équation \deleted{homogène}
\[
 \sum x_i=0 . \tag*{(13)}
\]
% slanted
\margin{description valable loc. ét. -- la description intrinsèque est \deleted{\unsure{par}} l'hyperplan (13) de $\check{\mathbb P}(p_*(\mathcal O_{\underline J}))$, d'équation $L\deleted{\unsure{(u)}}=0$, où $L(u)=\mathrm{Tr}\,u$}

\deleted{Donc} La donnée de $(X,\underline I)$ équivaut à celle de \deleted{$\unsure{\Sigma}$} $\underline J$ i.e de $\Sigma_{\underline J}$, du torseur $\underline K$ sous $V_{\underline J}$ i.e. du revêtement \deleted{\unsure{en}} $X$ de $\Sigma_{\underline J}$ \deleted{ramifié} de groupe $V_{\underline J}$ (admissible au sens qu'il définit l'\uline{identité} $V_{\underline J}\longrightarrow V_{\underline J}$), \uline{et} d'une section $\unsure{t}$ de $\Sigma_{\underline J}^*$.
% slanted
\margin{Relation entre $\underline I,\underline K,x$~: $(14)\;\big|\;\underline I\simeq\underline K\mathbin{\wedge}_{\underline V}(\overline\Sigma^*_{\underline J,2}|_{t})$. Ainsi, pour $\underline J$ et $x\in\Gamma(\Sigma^*_{\underline J}/S)$ fixés, $\underline I$ et $\underline K$ (ou encore, $\underline I$ et $X$) se déterminent mutuellement \ldots{} On se gardera de confondre les trois rev\textsuperscript{ts} étales de degré 4 de $S$, $\underline I$, $\underline J\amalg S$, et $\underline K$ \ldots}

Quand $\underline I$ est \inserted{fixé et} de la forme $I_S$, \inserted{[l'hypothèse que $\underline I$ soit constant est en fait superflue]} avec $I$ ens. à quatre éléments fixé (à fins d'épinglage), on a $\underline J=J_S$, et $\Sigma_{\underline J}^*=\overline\Sigma_J$ est aussi fixé. La donnée de $(X,I_S\hookrightarrow X)$ équivaut alors à celle de la section $x$ de $\Sigma_J^*$, servant de ``pt référence'', le revêtement $X$ de $\Sigma_J^*$ de groupe $(V_J)_S=(V)_{(J_S)}=\underline V$ étant
